\documentclass{article}
\usepackage[preprint]{neurips_2023}

\usepackage[utf8]{inputenc}
\usepackage[T1]{fontenc}
\usepackage{hyperref}
\usepackage{url}
\usepackage{booktabs}
\usepackage{amsfonts}
\usepackage{nicefrac}
\usepackage{microtype}
\usepackage{xcolor}
\usepackage{amsmath}

\title{Name That LLM: Fine-Grained Attribution of Short Texts}

\author{%
  Mat\'ias Manuel Zeballos \\
  ACP Master Student\\
  \texttt{2026280067} \\
  \And
  Helina Mesfin Geresu \\
  ACP Master Student\\
  \texttt{2026280061} \\
  \And
  Ma\"eli Carcass\`es \\
  ACP Master Student\\
  \texttt{2026280073} \\
}

\begin{document}
\maketitle

\begin{abstract}
Identifying which language model produced a short passage is a fine-grained authorship attribution problem. Unlike binary AI-text detection, it requires distinguishing among models that all generate machine-written text. We frame our group project around Kaggle's \textit{Name That LLM} competition, which asks participants to identify the authoring model for passages of varying length across 18 models from six AI labs. We propose comparing word-level and character-level TF--IDF classifiers with a pretrained text encoder, using validation procedures and balanced accuracy that reflect the competition's class-imbalanced evaluation. The project notes currently record the team and task but contain no experimental results; this paper therefore reports a study plan rather than performance claims.
\end{abstract}

\section{Introduction}
Text generated by large language models is increasingly common, making it useful to know not only whether a passage is machine-written, but which model produced it. Model attribution may support analysis of synthetic text provenance and model-specific stylistic differences. The \textit{Name That LLM} competition poses this problem as an 18-class classification task over short, self-contained responses to prompts from several subject areas \citep{yan2026namethatllm}.

Attribution is challenging because passages are short and vary in length, prompts are not necessarily seen during evaluation, and models from the same organization may share stylistic characteristics. A useful system should therefore capture linguistic patterns beyond topic words, and evaluation should give each candidate model equal importance.

\section{Task and evaluation}
The goal is to assign each unlabeled passage to exactly one of 18 candidate language models. The competition describes test passages of approximately 40--250 words, with roughly 4,000 examples in the hidden test set. The test set's class distribution is not uniform, and its model mix may differ from the public leaderboard. Accordingly, the competition uses balanced accuracy:
\begin{equation}
  \operatorname{BA} = \frac{1}{K} \sum_{k=1}^{K}
  \frac{\operatorname{TP}_k}{\operatorname{TP}_k + \operatorname{FN}_k},
  \qquad K=18.
\end{equation}
This is the mean recall over the candidate models, so a strong score requires performance across classes rather than success on only the most frequent labels.

\section{Proposed methodology}
\paragraph{Data inspection.} We will inspect the available training data for class frequencies, passage lengths, missing text, and repeated or near-duplicate passages. Any preprocessing will be fitted using training data only. We will avoid removing punctuation or normalizing casing by default, since these choices could discard stylistic signals relevant to attribution.

\paragraph{Baselines.} We will establish a majority-class baseline and a word-level TF--IDF model with a linear classifier. The word baseline provides a simple reference point, while a character n-gram TF--IDF classifier can capture spelling, punctuation, and subword habits. A combined word-and-character feature model will test whether these complementary signals improve validation performance. Hyperparameters will be selected using validation data, with balanced accuracy as the primary metric.

\paragraph{Neural comparison.} If the available data and compute budget permit, we will compare the feature-based systems with a pretrained text encoder and a classification head. We will tune regularization and learning rate on validation data and use early stopping. This comparison will help determine whether contextual representations add value over compact lexical and character-level signals; it is a proposed experiment, not a completed result.

\section{Experimental protocol}
We will reserve a stratified validation partition from the labeled training data, preserving examples from each class in both partitions. If prompt identifiers are available, we will also consider a prompt-disjoint validation split to better approximate evaluation on unseen prompts. We will report overall balanced accuracy and per-class recall, and inspect the confusion matrix to identify model pairs that are difficult to separate. The final competition submission will contain one valid model label for every test identifier, in the required format.

\section{Limitations and current status}
The project notes supplied for this draft contain the team roster and competition name, but no dataset analysis, implementation details, validation scores, or submission results. Consequently, the methods above are planned experiments and no claim of achieved performance is made. Results, implementation choices, and error analysis should be added after experiments are run. The competition's hidden test distribution may also differ from any local validation split, so validation performance alone cannot establish leaderboard performance.

\section*{Acknowledgments}
This draft is based on the team notes and the public competition description. No funding or experimental results were reported in the supplied notes.

\begin{thebibliography}{1}
\small
\bibitem[Yan et~al.(2026)Yan, Plomecka, and Oldacre]{yan2026namethatllm}
Yao Yan, Martyna Plomecka, and Ashley Oldacre.
\newblock Name That LLM.
\newblock Kaggle competition, 2026.
\newblock \url{https://www.kaggle.com/competitions/name-that-llm}.
\end{thebibliography}

\end{document}
